
\section{\label{sec:Extrapolation}Ekstrapolasi}

Dalam kalkulus, Anda tentu pernah menjumpai konstanta Euler, $e$,
yang mungkin diberitahukan kepada Anda bernilai kira-kira $2.718$,
atau mungkin hanya $2.7$. Dan kecuali jika Anda pernah mengikuti
lomba menghafal digit-digit $e$, Anda mungkin tidak pernah melihat
lebih banyak digit $e$ daripada yang dapat ditampilkan kalkulator
Anda. Kita akan segera mengubahnya. Lima puluh digit pertama $e$
adalah
\[
2.7182818284590452353602874713526624977572470936999.
\]
Berapa banyak yang Anda ingat? Tidak perlu khawatir jika tidak terlalu
banyak. Tidak akan ada kuis tentang digit-digit $e$ dalam waktu dekat. 

\noindent\begin{digression}{Digit-digit $e$} 

\noindent Seribu digit pertama $e$, sebanyak 50 digit per baris,
adalah
\[
\begin{gathered}2.7182818284590452353602874713526624977572470936999\\
59574966967627724076630353547594571382178525166427\\
42746639193200305992181741359662904357290033429526\\
05956307381323286279434907632338298807531952510190\\
11573834187930702154089149934884167509244761460668\\
08226480016847741185374234544243710753907774499206\\
95517027618386062613313845830007520449338265602976\\
06737113200709328709127443747047230696977209310141\\
69283681902551510865746377211125238978442505695369\\
67707854499699679468644549059879316368892300987931\\
27736178215424999229576351482208269895193668033182\\
52886939849646510582093923982948879332036250944311\\
73012381970684161403970198376793206832823764648042\\
95311802328782509819455815301756717361332069811250\\
99618188159304169035159888851934580727386673858942\\
28792284998920868058257492796104841984443634632449\\
68487560233624827041978623209002160990235304369941\\
84914631409343173814364054625315209618369088870701\\
67683964243781405927145635490613031072085103837505\\
10115747704171898610687396965521267154688957035035
\end{gathered}
\]

\noindent\end{digression}

\noindent Namun, ingatkah Anda dari kalkulus bahwa 
\[
\lim_{h\to0}(1+h)^{1/h}=e?
\]
Dapatkah Anda membuktikannya? Bukti \vpageref{proofOfLimit}. Berdasarkan
fakta ini, kita dapat menggunakan 
\[
\tilde{e}(h)=(1+h)^{1/h}
\]
untuk menghampiri $e$. Mari langsung mencobanya!
\begin{eqnarray*}
\tilde{e}(0.01) & \approx & 2.704813829421529\\
\tilde{e}(0.005) & \approx & 2.711517122929293\\
\tilde{e}(0.0025) & \approx & 2.714891744381238\\
\tilde{e}(0.00125) & \approx & 2.716584846682473\\
\tilde{e}(0.000625) & \approx & 2.717432851769196.
\end{eqnarray*}
Sayangnya, barisan hampiran ini tidak konvergen dengan cepat. Hampiran
pertama memiliki dua digit akurat, sedangkan hampiran kelima pun masih
hanya memiliki tiga digit akurat. Tentu saja, kita dapat terus memperkecil
$h$ untuk memperoleh hampiran yang lebih akurat, tetapi berdasarkan
lambatnya peningkatan sejauh ini, cara tersebut tampaknya kurang menjanjikan.
Sebagai gantinya, kita dapat menggabungkan taksiran yang sudah ada
untuk memperoleh hampiran yang lebih baik. Gagasan ini, setidaknya
secara sepintas, seharusnya mengingatkan Anda pada metode delta-kuadrat
Aitken. Dalam metode itu, tiga hampiran berurutan digabungkan untuk
membentuk hampiran baru yang umumnya lebih baik daripada ketiga hampiran
asalnya. Di sini kita akan melakukan hal serupa: menggabungkan hampiran
yang kurang memadai untuk memperoleh hampiran yang lebih baik. Berbagai
hampiran baru ini akan kita beri nama agar dapat terus digunakan.
\begin{eqnarray}
2\tilde{e}(0.005)-\tilde{e}(0.01) & \equiv & \tilde{e}_{1}(0.01)=2.718220416437056\nonumber \\
2\tilde{e}(0.0025)-\tilde{e}(0.005) & \equiv & \tilde{e}_{1}(0.005)=2.718266365833184\nonumber \\
2\tilde{e}(0.00125)-\tilde{e}(0.0025) & \equiv & \tilde{e}_{1}(0.0025)=2.718277948983707\nonumber \\
2\tilde{e}(0.000625)-\tilde{e}(0.00125) & \equiv & \tilde{e}_{1}(0.00125)=2.718280856855920.\label{eq:richardsonC2}
\end{eqnarray}
Setiap hampiran baru ini akurat hingga 5 atau 6 digit signifikan!
Itu sudah merupakan peningkatan besar. Kita dapat menggabungkannya
lebih lanjut untuk memperoleh hampiran yang lebih baik lagi:
\begin{eqnarray}
\frac{4\tilde{e}_{1}(0.005)-\tilde{e}_{1}(0.01)}{3} & \equiv & \tilde{e}_{2}(0.01)=2.718281682298560\nonumber \\
\frac{4\tilde{e}_{1}(0.0025)-\tilde{e}_{1}(0.005)}{3} & \equiv & \tilde{e}_{2}(0.005)=2.718281810033881\nonumber \\
\frac{4\tilde{e}_{1}(0.00125)-\tilde{e}(0.0025)}{3} & \equiv & \tilde{e}_{2}(0.0025)=2.718281826146657.\label{eq:richardsonC3}
\end{eqnarray}
Hampiran pertama akurat hingga tujuh digit signifikan, yang kedua
hingga delapan, dan yang ketiga hingga sembilan! Kita masih dapat
menggabungkannya lebih lanjut:
\begin{eqnarray}
\frac{8\tilde{e}_{2}(0.005)-\tilde{e}_{2}(0.01)}{7} & \equiv & \tilde{e}_{3}(0.01)=2.718281828281785\nonumber \\
\frac{8\tilde{e}_{2}(0.0025)-\tilde{e}_{2}(0.005)}{7} & \equiv & \tilde{e}_{3}(0.005)=2.718281828448482.\label{eq:richardsonC4}
\end{eqnarray}
Kini kita memiliki hampiran yang akurat hingga sepuluh dan sebelas
digit signifikan! Jika ditinjau kembali, kita mengambil lima hampiran
yang masing-masing akurat tidak lebih dari 3 digit signifikan, lalu
menggabungkannya menjadi dua hampiran yang masing-masing akurat hingga
sedikitnya 10 digit signifikan. Ajaib! Baiklah, bukan sulap, melainkan
keajaiban matematika! Berikut cara kerjanya.

Misalkan kita menghampiri $p$ dengan menggunakan rumus $\tilde{p}(h)$,
dan kita mengetahui bahwa
\[
\tilde{p}(h)=p+c_{1}\cdot h^{m_{1}}+c_{2}\cdot h^{m_{2}}+c_{3}\cdot h^{m_{3}}+\cdots.
\]
Maka
\[
\tilde{p}\left(\alpha h\right)=p+c_{1}\cdot\left(\alpha h\right)^{m_{1}}+c_{2}\cdot(\alpha h)^{m_{2}}+c_{3}\cdot(\alpha h)^{m_{3}}+\cdots.
\]
Sekarang, jika persamaan kedua dikalikan dengan $\alpha^{-m_{1}}$
lalu persamaan pertama dikurangkan darinya, suku $h^{m_{1}}$ lenyap,
dan kita memperoleh hampiran dengan suku galat yang diawali oleh $c_{2}\cdot h^{m_{2}}$:

\begin{center}
\begin{tabular}{rcl}
$\alpha^{-m_{1}}\tilde{p}\left(\alpha h\right)$ & $=$ & $\alpha^{-m_{1}}p+c_{1}\cdot h^{m_{1}}+c_{2}\alpha^{m_{2}-m_{1}}\cdot h^{m_{2}}+c_{3}\alpha^{m_{3}-m_{1}}\cdot h^{m_{3}}+\cdots$\tabularnewline
$-\left[\tilde{p}(h)\right.$ & $=$ & $\left.p+c_{1}\cdot h^{m_{1}}+c_{2}\cdot h^{m_{2}}+c_{3}\cdot h^{m_{3}}+\cdots\right]$\tabularnewline
\hline 
$\alpha^{-m_{1}}\tilde{p}\left(\alpha h\right)-\tilde{p}(h)$ & $=$ & $(\alpha^{-m}-1)p+c_{2}(\alpha^{m_{2}-m_{1}}-1)\cdot h^{m_{2}}+c_{3}(\alpha^{m_{3}-m_{1}}-1)\cdot h^{m_{3}}+\cdots$\tabularnewline
\end{tabular}
\par\end{center}

\noindent Dengan sedikit penataan ulang,
\begin{equation}
\frac{\alpha^{-m_{1}}\tilde{p}\left(\alpha h\right)-\tilde{p}(h)}{\alpha^{-m_{1}}-1}=p+d_{2}\cdot h^{m_{2}}+d_{3}\cdot h^{m_{3}}+\cdots\label{eq:richardsonBasis}
\end{equation}
untuk beberapa konstanta $d_{2},d_{3},\ldots$. Jika $m_{2}>m_{1}$,
metode ini cenderung memperbaiki kedua hampiran $\tilde{p}(h)$ dan
$\tilde{p}\left(\alpha h\right)$ dengan menggabungkannya menjadi
satu hampiran yang galatnya seorde dengan suatu kelipatan konstan
dari $h^{m_{2}}$. Perhitungan ini menjadi dasar ekstrapolasi Richardson.\index{ekstrapolasi Richardson}

Secara kebetulan, $\tilde{e}(h)$ memiliki tepat bentuk yang diperlukan.
\begin{equation}
\tilde{e}(h)=e+c_{1}h+c_{2}h^{2}+c_{3}h^{3}+c_{4}h^{4}+O(h^{5})\label{eq:eErrorForm}
\end{equation}
untuk beberapa konstanta $c_{1},c_{2},c_{3},c_{4}$. Nilai sebenarnya
dari konstanta-konstanta tersebut tidak relevan bagi perhitungan ini.
Untuk memahami perhitungan $\tilde{e}_{1},$, kita menggunakan Persamaan
\ref{eq:richardsonBasis} dengan $\alpha=\frac{1}{2}$ dan $m_{1}=1$
untuk memperoleh
\begin{eqnarray*}
\tilde{e}_{1}(h) & = & \frac{2\tilde{e}\left(\frac{h}{2}\right)-\tilde{e}(h)}{2-1}\\
 & = & 2e+c_{1}h+\frac{1}{2}c_{2}h^{2}+\frac{1}{4}c_{3}h^{3}+\frac{1}{8}c_{4}h^{4}+O(h^{5})\\
 &  & \qquad-\left[e+c_{1}h+c_{2}h^{2}+c_{3}h^{3}+c_{4}h^{4}+O(h^{5})\right]\\
 & = & e+d_{2}h^{2}+d_{3}h^{3}+d_{4}h^{4}+O(h^{5})
\end{eqnarray*}
untuk beberapa konstanta $d_{2},d_{3},d_{4}$. $\tilde{e}_{1}(h)$
merupakan rumus yang menghasilkan deretan hampiran akurat hingga 5
atau 6 digit signifikan. Menentukan konstanta $d_{i}$ berdasarkan
konstanta $c_{i}$ tidaklah sulit, tetapi sekali lagi nilai konstanta-konstanta
itu tidak penting dan hanya akan memperumit penyempurnaan selanjutnya.
Yang penting adalah bentuk galatnya. Setelah mengetahui $\tilde{e}_{1}(h)=e+d_{2}h^{2}+d_{3}h^{3}+d_{4}h^{4}+O(h^{5})$,
kita memperoleh $\tilde{e}_{2}(h)$ dengan menggunakan rumus \ref{eq:richardsonBasis}
dengan $\alpha=\frac{1}{2}$ dan $m_{1}=2$:
\begin{eqnarray*}
\tilde{e}_{2}(h) & = & \frac{4\tilde{e}_{1}\left(\frac{h}{2}\right)-\tilde{e}_{1}(h)}{3}\\
 & = & e+k_{3}h^{3}+k_{4}h^{4}+O(h^{5})
\end{eqnarray*}
untuk beberapa konstanta $k_{3}$ dan $k_{4}$. $\tilde{e}_{2}(h)$
merupakan rumus yang menghasilkan deretan hampiran akurat hingga 7
sampai 9 digit signifikan. Kita kembali dapat menggunakan rumus \ref{eq:richardsonBasis},
kali ini dengan $\alpha=\frac{1}{2}$ dan $m_{1}=3$:
\begin{eqnarray*}
\tilde{e}_{3}(h) & = & \frac{8\tilde{e}_{2}\left(\frac{h}{2}\right)-\tilde{e}_{2}(h)}{7}\\
 & = & e+l_{4}h^{4}+O(h^{5})
\end{eqnarray*}
untuk suatu konstanta $l_{4}$. $\tilde{e}_{3}(h)$ merupakan rumus
yang menghasilkan hampiran akurat hingga 10 dan 11 digit signifikan.
Sekarang cobalah menggunakan bentuk $\tilde{e}_{3}(h)$ dan rumus
\ref{eq:richardsonBasis} untuk menurunkan rumus $O(h^{5})$ bagi
$\tilde{e}_{4}(h)$. Kemudian gunakan rumus Anda untuk menghitung
$\tilde{e}_{4}(0.01)$ dengan menggunakan nilai $\tilde{e}_{3}(0.01)$
dan $\tilde{e}_{3}(0.005)$. Seberapa akurat $\tilde{e}_{4}(0.01)?$
Jawaban \vpageref{eTilde4}.

Sebagai kasus khusus, penerapan ekstrapolasi Richardson dengan $\alpha=\frac{1}{2}$
pada sebarang hampiran berbentuk
\[
\tilde{p}_{0}(h)=p+c_{1}h+c_{2}h^{2}+c_{3}h^{3}+\cdots
\]
menghasilkan penyempurnaan yang didefinisikan secara rekursif
\[
\tilde{p}_{k}(h)=\frac{2^{k}\tilde{p}_{k-1}\left(\frac{h}{2}\right)-\tilde{p}_{k-1}(h)}{2^{k}-1},\quad k=1,2,3,\ldots
\]
yang diharapkan semakin akurat ketika $k$ bertambah. Untuk nilai
$\alpha$ lain atau bentuk galat lain, rumus bagi $\tilde{p}_{k}(h)$
berubah sesuai dengan \ref{eq:richardsonBasis}.

\noindent\begin{digression}{Polinom Taylor untuk $\tilde e(h)$} \label{crumpet:eApproximate}

\noindent$\tilde{e}$ tidak terdefinisi pada $0$, sehingga turunannya
pada $0$ juga tidak terdefinisi. Namun, jika kita memperluas definisi
$\tilde{e}$ menjadi
\[
\tilde{e}(h)=\begin{cases}
(1+h)^{1/h} & \mbox{ jika }h\neq0\\
e & \mbox{ jika }h=0
\end{cases},
\]
, dengan demikian mendefinisikan $\tilde{e}$ pada $0$, maka $\tilde{e}(h)$
dapat didiferensialkan tak berhingga kali pada $0$, dan polinom Taylor
kelimanya, misalnya, adalah:

\noindent
\[
\tilde{e}(h)=e-\frac{e}{2}\cdot h+\frac{11e}{24}\cdot h^{2}-\frac{7e}{16}\cdot h^{3}+\frac{2447e}{5760}\cdot h^{4}+\frac{f^{(5)}(\xi)}{120}h^{5}
\]
untuk suatu $\xi\in(0,h)$.

\noindent\end{digression} 

\subsection{Diferensiasi}

Dengan ekstrapolasi, rumus hampiran diferensiasi berorde tinggi dapat
diturunkan dari rumus berorde rendah. Kita mulai dengan hampiran berorde
terendah, ${\displaystyle f'(x_{0})=\frac{-f(x_{0})+f(x_{0}+h)}{h}-\frac{h}{2}f''(\xi_{h})}$.
Suku galat standar, $-\frac{h}{2}f''(\xi_{h})$ tidak menyatakan galat
dalam bentuk $c\cdot h^{m_{1}}+O(h^{m_{2}})$ sebagaimana disyaratkan
oleh ekstrapolasi Richardson, sehingga kita kembali ke deret Taylor
untuk menentukan suku $O(h^{m_{2}})$ berikut:
\[
f(x_{0}+h)=f(x_{0})+hf'(x_{0})+\frac{1}{2}h^{2}f''(x_{0})+\frac{1}{6}h^{3}f'''(x_{0})+\cdots
\]
sehingga
\[
\frac{-f(x_{0})+f(x_{0}+h)}{h}=f'(x_{0})+\frac{1}{2}hf''(x_{0})+\frac{1}{6}h^{2}f'''(x_{0})+\cdots.
\]
Dengan demikian, 
\begin{eqnarray*}
f'(x_{0})-\frac{-f(x_{0})+f(x_{0}+h)}{h} & = & -\frac{1}{2}hf''(x_{0})-\frac{1}{6}h^{2}f'''(x_{0})-\cdots\\
 & = & c_{1}h+O(h^{2})
\end{eqnarray*}
dan ekstrapolasi akan menghasilkan rumus dengan galat $O(h^{2})$.
Dengan menetapkan $\tilde{p}(h)=\frac{-f(x_{0})+f(x_{0}+h)}{h}$,
$\alpha=2$, dan $m_{1}=1$, rumus \ref{eq:richardsonBasis} memberikan
hampiran
\[
\frac{\frac{1}{2}\tilde{p}(2h)-\tilde{p}(h)}{\frac{1}{2}-1}
\]
yang akan menjadi rumus dengan galat $O(h^{2})$ untuk $f'(x_{0})$.
Setelah disederhanakan,
\begin{eqnarray*}
\frac{\frac{1}{2}\tilde{p}(2h)-\tilde{p}(h)}{\frac{1}{2}-1} & = & \frac{\frac{1}{2}\left[\frac{-f(x_{0})+f(x_{0}+2h)}{2h}\right]-\frac{-f(x_{0})+f(x_{0}+h)}{h}}{-\frac{1}{2}}\\
 & = & \frac{\frac{-f(x_{0})+f(x_{0}+2h)}{4h}-\frac{-4f(x_{0})+4f(x_{0}+h)}{4h}}{-\frac{1}{2}}\\
 & = & \frac{\frac{3f(x_{0})-4f(x_{0}+h)+f(x_{0}+2h)}{4h}}{-\frac{1}{2}}\\
 & = & \frac{-3f(x_{0})+4f(x_{0}+h)-f(x_{0}+2h)}{2h}.
\end{eqnarray*}
Dengan demikian, kita memperoleh $f'(x_{0})=\frac{-3f(x_{0})+4f(x_{0}+h)-f(x_{0}+2h)}{2h}+O(h^{2})$,
tetapi hasil ini bukan hal baru. Ini adalah rumus 3 titik pertama
dalam Tabel \ref{tab:firstDerivativeFormulas}! Rumus turunan berorde
tinggi lainnya juga dapat diturunkan dengan ekstrapolasi, tetapi umumnya
hasilnya tidak memberi pengetahuan baru. Kita hanya memperoleh cara
baru untuk menurunkan rumus diferensiasi berorde tinggi.

\subsection{Integrasi}

Penerapan ekstrapolasi pada integral tentu memberikan hasil yang lebih
bermanfaat. Kita mulai dengan sebarang rumus integrasi komposit dan
menerapkan ekstrapolasi Richardson. Sekarang kita tinjau aturan trapesium
komposit dan menggunakan notasi $T_{k}(a,b)$ untuk menyatakan hampiran
$\int_{a}^{b}f(x)dx$ dengan aturan trapesium yang menggunakan $k$
subselang.

Sebelum melanjutkan, kita perlu memahami dengan baik arti aturan trapesium
komposit yang memiliki suku galat $O\left(\left(\frac{1}{n}\right)^{2}\right)$.
Pada dasarnya, kita mengharapkan galat berkurang dengan faktor sekitar
4 ketika jumlah interval digandakan. Kita mengharapkan galat berkurang
dengan faktor sekitar 9 ketika jumlah interval dibuat tiga kali lipat.
Secara umum, kita mengharapkan galat berkurang dengan faktor sekitar
$\beta^{2}$ ketika jumlah interval dikalikan dengan $\beta$. Untuk
melihat dampak ini secara nyata, tinjau integral tentu
\[
\int_{0}^{1}\sin x\,dx
\]
yang nilai eksaknya adalah $1-\cos(1)\approx.4596976941318602$. Galat
mutlak dari $T_{5}(0,1)$, $T_{10}(0,1),$ dan $T_{15}(0,1)$ adalah
\begin{eqnarray*}
\left|\int_{0}^{1}\sin x\,dx-T_{5}(0,1)\right| & \approx & 1.533(10)^{-3}\\
\left|\int_{0}^{1}\sin x\,dx-T_{10}(0,1)\right| & \approx & 3.831(10)^{-4}\\
\left|\int_{0}^{1}\sin x\,dx-T_{15}(0,1)\right| & \approx & 1.702(10)^{-4}
\end{eqnarray*}
Kita mengharapkan galat $\left|\int_{0}^{1}\sin x\,dx-T_{5}(0,1)\right|$
sekitar empat kali galat $\left|\int_{0}^{1}\sin x\,dx-T_{10}(0,1)\right|$
dan sembilan kali galat $\left|\int_{0}^{1}\sin x\,dx-T_{15}(0,1)\right|$.
Untuk memeriksanya, kita hitung rasionya:
\begin{eqnarray*}
\frac{\left|\int_{0}^{1}\sin x\,dx-T_{5}(0,1)\right|}{\left|\int_{0}^{1}\sin x\,dx-T_{10}(0,1)\right|} & = & \frac{1.533(10)^{-3}}{3.831(10)^{-4}}\approx4.001\\
\frac{\left|\int_{0}^{1}\sin x\,dx-T_{5}(0,1)\right|}{\left|\int_{0}^{1}\sin x\,dx-T_{15}(0,1)\right|} & = & \frac{1.533(10)^{-3}}{1.702(10)^{-4}}\approx9.007.
\end{eqnarray*}
Kira-kira berapa nilai rasio $\frac{\left|\int_{0}^{1}\sin x\,dx-T_{10}(0,1)\right|}{\left|\int_{0}^{1}\sin x\,dx-T_{15}(0,1)\right|}$
yang Anda harapkan? Jawaban \vpageref{errorRatio}.

Akhirnya, kita menerapkan ekstrapolasi Richardson dengan $\alpha=\frac{1}{2}$
dan $m_{1}=2$ untuk menghasilkan taksiran berorde lebih tinggi,
\[
T_{k,1}(a,b)\equiv\frac{4T_{2k}(a,b)-T_{k}(a,b)}{3}.
\]
Kita menggunakan perhitungan numerik untuk memahami suku galat dari
penyempurnaan $T_{k,1}$. Kita mulai dengan mengumpulkan beberapa
data. Dengan melanjutkan analisis $\int_{0}^{1}\sin x\,dx$, perhatikan
bahwa
\begin{eqnarray*}
T_{5}(0,1) & \approx & .4581643459604436\\
T_{10}(0,1) & \approx & .4593145488579763\\
T_{20}(0,1) & \approx & .4596019197882473\\
T_{40}(0,1) & \approx & .4596737512942187.
\end{eqnarray*}
Dengan demikian,
\begin{eqnarray*}
T_{5,1}(0,1) & = & \frac{4T_{10}(0,1)-T_{5}(0,1)}{3}\approx.4596979498238206\\
T_{10,1}(0,1) & = & \frac{4T_{20}(0,1)-T_{10}(0,1)}{3}\approx.4596977100983375\\
T_{20,1}(0,1) & = & \frac{4T_{40}(0,1)-T_{20}(0,1)}{3}\approx.4596976951295424
\end{eqnarray*}
dan
\begin{eqnarray*}
\frac{\left|\int_{0}^{1}\sin x\,dx-T_{5,1}(0,1)\right|}{\left|\int_{0}^{1}\sin x\,dx-T_{10,1}(0,1)\right|} & \approx & 16.01\\
\frac{\left|\int_{0}^{1}\sin x\,dx-T_{10,1}(0,1)\right|}{\left|\int_{0}^{1}\sin x\,dx-T_{20,1}(0,1)\right|} & \approx & 16.00.
\end{eqnarray*}
Ketika jumlah subselang digandakan, galat berkurang dengan faktor
16. Nilai itu adalah $2^{4}$, bukan $2^{3}$ seperti yang mungkin
kita duga! Penyempurnaan pertama meningkatkan hampiran dengan galat
$O\left(\left(\frac{1}{n}\right)^{2}\right)$ menjadi hampiran dengan
galat $O\left(\left(\frac{1}{n}\right)^{4}\right)$. Dengan kata lain,
galat $T_{n,1}$ adalah $O\left(\left(\frac{1}{n}\right)^{4}\right)$.

Setelah mengetahui bahwa galat $T_{n,1}$ adalah $O\left(\left(\frac{1}{n}\right)^{4}\right)$,
kita dapat melakukan ekstrapolasi lagi. Dengan menerapkan ekstrapolasi
Richardson dengan $\alpha=\frac{1}{2}$ dan $m_{1}=4$, kita memperoleh
\begin{eqnarray*}
T_{5,2}(0,1) & = & \frac{16T_{10,1}(0,1)-T_{5,1}(0,1)}{15}\approx.4596976941166387\\
T_{10,2}(0,1) & = & \frac{16T_{20,1}(0,1)-T_{10,1}(0,1)}{15}\approx.4596976941316228.
\end{eqnarray*}
Kini kita memiliki hampiran $T_{5,2}$ dan $T_{10,2}$ yang galatnya
masing-masing hanya sekitar $1.522(10)^{-11}$ dan $2.374(10)^{-13}$.
Gunakan informasi ini untuk menghitung $T_{5,3}$ beserta galat mutlaknya.
Jawaban \vpageref{t35}.

Metode yang menggabungkan ekstrapolasi Richardson dengan aturan trapesium
dikenal sebagai metode Romberg atau integrasi Romberg.\index{integrasi numerik!Romberg}\index{integrasi Romberg}
Perhitungan sering ditabelkan untuk memudahkan penataan, seperti dalam
Tabel 
\begin{table}
\hrule \caption{\label{tab:Romberg}Metode Romberg}

\centering{}\medskip{}
\begin{tabular}{ccccc}
$T_{1}$ & $T_{1,1}$ & $T_{1,2}$ & $T_{1,3}$ & $\cdots$\tabularnewline
$T_{2}$ & $T_{2,1}$ & $T_{2,2}$ & $\vdots$ & \tabularnewline
$T_{4}$ & $T_{4,1}$ & $\vdots$ &  & \tabularnewline
$T_{8}$ & $\vdots$ &  &  & \tabularnewline
$\vdots$ &  &  &  & \tabularnewline
\end{tabular}\medskip{}
\hrule 
\end{table}
 \ref{tab:Romberg}. Baris ditambahkan sampai kedua selisih $\left|T_{k,n}-T_{k,n+1}\right|$
dan $\left|T_{2k,n}-T_{k,n+1}\right|$ lebih kecil daripada suatu
toleransi.

Meskipun ekstrapolasi Richardson dapat diterapkan pada sebarang rumus
integrasi komposit, perhitungan suku galat di atas membantu menjelaskan
mengapa aturan trapesium merupakan pilihan yang tepat. Dari perhitungan
tersebut, kita dapat menduga (dan hal ini dapat dibuktikan) bahwa
suku galat aturan trapesium komposit hanya memuat pangkat genap dari
$\frac{1}{n}$. Secara eksplisit, kita memperoleh
\[
\int_{a}^{b}f(x)dx=T_{n}(a,b)+c_{2}\left(\frac{1}{n}\right)^{2}+c_{4}\left(\frac{1}{n}\right)^{4}+c_{6}\left(\frac{1}{n}\right)^{6}+\cdots
\]
sehingga setiap penyempurnaan menaikkan derajat terendah dalam suku
galat sebesar 2, bukan 1. Dengan melewati derajat ganjil, pilihan
ini menjadi sangat efisien. Namun, metode ini memiliki konsekuensi.
Di balik $c_{2}$ tersimpan asumsi bahwa $f$ memiliki turunan kedua
yang kontinu. Di balik $c_{4}$ tersimpan asumsi bahwa $f$ memiliki
turunan keempat yang kontinu. Demikian seterusnya. Keakuratan setiap
penyempurnaan mensyaratkan $f$ yang memiliki dua turunan kontinu
tambahan. Semakin banyak penyempurnaan yang dilakukan, $f$ harus
semakin mulus agar metode ini bekerja. Karena itu, metode Romberg
sebaiknya digunakan hanya apabila integran diketahui memiliki turunan
yang memadai.

\subsection{Konsep Utama}
\begin{description}
\item [{Ekstrapolasi~Richardson:}] Jika hampiran $\tilde{p}$ diketahui
berbentuk
\[
\tilde{p}(h)=p+c_{1}h^{m_{1}}+O(h^{m_{2}})
\]
maka hampiran 
\[
\frac{\alpha^{-m_{1}}\tilde{p}\left(\alpha h\right)-\tilde{p}(h)}{\alpha^{-m_{1}}-1}
\]
akan memiliki galat $O(h^{m_{2}})$.
\item [{Integrasi~Romberg:}] Penerapan ekstrapolasi Richardson pada metode
trapesium.
\end{description}
\noindent\startexercises 

\subsection{Latihan}
\begin{enumerate}
\item \label{exc:calculus-extrapolation}Polinom Taylor dapat digunakan
untuk menunjukkan bahwa
\[
\pi=2\sin^{-1}(1-h^{4})+K_{2}h^{2}+K_{4}h^{6}+K_{6}h^{10}+\cdots.
\]
Karena itu, $N(h)=2\sin^{-1}(1-h^{4})$ merupakan hampiran dengan
galat $O(h^{2})$ untuk $\pi$. Gunakan ekstrapolasi Richardson untuk
menurunkan hampiran dengan galat $O(h^{6})$ untuk $\pi$. \hasananswer 
\item Menarik untuk diperhatikan bahwa kita dapat merekayasa balik penyempurnaan
Richardson guna menghampiri $c_{i}$ pada persamaan \vref{eq:eErrorForm}.
Sebagai contoh, $\tilde{e}(h)=e+c_{1}h+O(h^{2})$, dan kita mengasumsikan
suku $O(h^{2})$ relatif kecil, sehingga persamaan ini dapat disusun
ulang untuk memperoleh
\[
\frac{\tilde{e}(h)-e}{h}\approx c_{1}.
\]
Sebagai contoh khusus, $\frac{\tilde{e}(.005)-e}{.005}=\frac{2.711517122929293-e}{.005}\approx-1.35$,
sehingga $c_{1}\approx-1.35$. Jika kita mencermati bagaimana konstanta-konstanta
tersebut berubah ketika kita menyempurnakan hampiran awal, kita juga
dapat memperoleh $c_{2}$, $c_{3}$, dan $c_{4}$. 
\begin{eqnarray*}
\tilde{e}_{1}(h) & = & 2\tilde{e}\left(\frac{h}{2}\right)-\tilde{e}(h)\\
 & = & 2e+c_{1}h+\frac{c_{2}}{2}h^{2}+\frac{c_{3}}{4}h^{3}+\frac{c_{4}}{8}h^{4}+O(h^{5})\\
 &  & -(e+c_{1}h+c_{2}h^{2}+c_{3}h^{3}+c_{4}h^{4}+O(h^{5}))\\
 & = & e-\frac{c_{2}}{2}h^{2}-\frac{3c_{3}}{4}h^{3}-\frac{7c_{4}}{8}h^{4}+O(h^{5}).
\end{eqnarray*}
Oleh karena itu, $\tilde{e}_{1}(h)-e\approx-\frac{c_{2}}{2}h^{2}$,
dan dari sini kita menyimpulkan 
\[
\frac{-2(\tilde{e}_{1}(h)-e)}{h^{2}}\approx c_{2}.
\]

\begin{enumerate}
\item Gunakan rumus ini dan nilai-nilai pada $\ref{eq:richardsonC2}$ untuk
memverifikasi bahwa $c_{2}\approx1.24$.
\item Hampiri $c_{3}$ menggunakan nilai-nilai pada \ref{eq:richardsonC3}.
\item Hampiri $c_{4}$ menggunakan nilai-nilai pada \ref{eq:richardsonC4}.
\item Bandingkan hampiran-hampiran untuk $c_{1},c_{2},c_{3},c_{4}$ ini
dengan nilai-nilai eksak pada kudapan \ref{crumpet:eApproximate}.
\end{enumerate}
\item \label{exc:calculus-extrapolation-1}Misalkan $N$ menghampiri $M$
menurut $N(h)=M+K_{1}h^{3}+K_{2}h^{5}+K_{3}h^{7}+\cdots$. Berorde
apakah $N_{3}(h)$ (ekstrapolasi Richardson generasi ketiga)? \hasananswer 
\item \label{exc:calculus-extrapolation-2}Misalkan $N$ menghampiri $M$
menurut $N(h)=M+K_{1}h^{2}+K_{2}h^{4}+K_{3}h^{6}+\cdots$. Kira-kira
berapakah nilai 
\[
\frac{\left|M-N(h/3)\right|}{\left|M-N(h/4)\right|}
\]
 untuk $h$ yang kecil? \hasananswer 
\item \label{exc:calculus-extrapolation-3}$N(h)=\frac{1-\cos h}{h^{2}}$
dapat digunakan untuk menghampiri \hasananswer  
\[
\lim_{h\to0}\frac{1-\cos h}{h^{2}}
\]

\begin{enumerate}
\item Hitung $N(1.0)$ dan $N(0.5)$.
\item \label{exc:calculus-extrapolation-4}Hitung $N_{1}(1.0)$, yaitu ekstrapolasi
Richardson pertama, dengan mengasumsikan

\begin{enumerate}
\item $N(h)$ memiliki galat berbentuk $K_{1}h+K_{2}h^{2}+K_{3}h^{3}+\cdots$
\item $N(h)$ memiliki galat berbentuk $K_{2}h^{2}+K_{4}h^{4}+K_{6}h^{6}+\cdots$
\end{enumerate}
\item Menurut Anda, asumsi manakah pada bagian \ref{exc:calculus-extrapolation-4}
yang memberikan bentuk galat yang benar, dan mengapa?
\end{enumerate}
\item Rumus beda mundur dapat dinyatakan sebagai
\[
\begin{gathered}f'(x_{0})=\frac{1}{h}[f(x_{0})-f(x_{0}-h)]\qquad\qquad\\
\qquad+\frac{h}{2}f''(x_{0})-\frac{h^{2}}{6}f'''(x_{0})+O(h^{3})
\end{gathered}
\]

\begin{enumerate}
\item Gunakan ekstrapolasi Richardson untuk menurunkan rumus $O(h^{2})$
bagi $f'(x_{0})$.
\item Rumus yang Anda turunkan seharusnya tampak familier. Rumus apakah
yang mirip dengannya? Apakah keduanya benar-benar sama? Mengapa atau
mengapa tidak?
\end{enumerate}
\item \label{exc:calculus-extrapolation-5}Turunkan rumus $O(h^{3})$ untuk
menghampiri $M$ yang menggunakan $N(h)$, $N(\frac{h}{2})$, dan
$N(\frac{h}{3})$, serta didasarkan pada asumsi bahwa \hasasolution 
\[
M=N(h)+K_{1}h+K_{2}h^{2}+K_{3}h^{3}+\cdots.
\]
\item \label{exc:calculus-extrapolation-6}Data berikut memberikan taksiran
integral $M=\int_{0}^{3\pi/2}\cos x\,dx$.
\[
\begin{gathered}N(h)=2.356194\qquad N(h/2)=-0.4879837\\
N(h/4)=-0.8815732\qquad N(h/8)=-0.9709157
\end{gathered}
\]
Dengan mengasumsikan $M-N(h)=K_{1}h^{2}+K_{2}h^{4}+K_{3}h^{6}+\cdots$,
tentukan ekstrapolasi Richardson ketiga untuk $M$. \hasasolution 
\item \label{exc:calculus-extrapolation-7}Misalkan $N(h)$ merupakan hampiran
bagi $M$ untuk setiap $h>0$ dan bahwa 
\[
M-N(h)=K_{1}h+K_{2}h^{2}+K_{3}h^{3}+\cdots
\]
untuk beberapa konstanta $K_{1},K_{2},K_{3},\ldots$. Gunakan nilai
$N(h)$, $N(h/3)$, dan $N(h/9)$ untuk menghasilkan hampiran $O(h^{3})$
bagi $M$. \hasananswer 
\item \label{exc:calculus-extrapolation-11}Gunakan integrasi Romberg untuk
menghitung integral dengan toleransi $10^{-4}$.

\begin{enumerate}
\item ${\displaystyle \int_{1}^{3}\ln(\sin(x))dx}$ \hasasolution 
\item ${\displaystyle \int_{5}^{7}\sqrt{x\cos x}\,dx}$
\item ${\displaystyle \int_{1}^{4}\frac{e^{x}\ln(x)}{x}dx}$ \hasananswer 
\item ${\displaystyle \int_{10}^{13}\sqrt{1+\cos^{2}x}\,dx}$
\item ${\displaystyle \int_{\ln3}^{\ln7}\frac{e^{x}}{1+x}dx}$ \hasananswer 
\item ${\displaystyle \int_{0}^{1}\frac{x^{2}-1}{x^{2}+1}dx}$
\item ${\displaystyle \int_{0}^{2}x^{2}\ln(x^{2}+1)dx}$ \hasananswer 
\end{enumerate}
\item \label{exc:calculus-extrapolation-8}\octave\ Tulislah fungsi Octave
yang menerapkan integrasi Romberg. \hasananswer 
\item \label{exc:calculus-extrapolation-10}\octave\ (i) Gunakan kode Anda
dari soal \ref{exc:calculus-extrapolation-8} untuk menghampiri integral
tersebut menggunakan $tol=10^{-5}$. (ii) Hitung galat sebenarnya
dari hampiran tersebut. (iii) Apakah hampiran tersebut memiliki galat
tidak lebih dari $10^{-5}$ seperti yang diminta?

\begin{enumerate}
\item \label{exc:calculus-extrapolation-9}${\displaystyle \int_{0}^{2\pi}x\sin(x^{2})dx}$ \hasananswer 
\item ${\displaystyle \int_{0.1}^{2}\frac{1}{x}\,dx}$
\item ${\displaystyle \int_{0}^{2}x^{2}\ln(x^{2}+1)\,dx}$
\end{enumerate}
CATATAN: $\int_{0}^{2}x^{2}\ln(x^{2}+1)\,dx=\frac{24\ln(5)-6\tan^{-1}(2)-4}{9}$.
\item Bandingkan hasil soal \ref{exc:calculus-extrapolation-10} dengan
hasil soal \vref{exc:calculus-composite-23}.
\end{enumerate}
\finishexercises 

\subsection{Jawaban}
\begin{description}
\item [{${\displaystyle \lim_{h\to0}(1+h)^{1/h}=e}$:}] \label{proofOfLimit}Mulailah
dengan memperhatikan bahwa $\ln\left[(1+h)^{1/h}\right]=\frac{\ln(1+h)}{h}$.
Tetapkan 
\begin{eqnarray*}
L & = & \lim_{h\to0}\frac{\ln(1+h)}{h}\\
 & = & \lim_{h\to0}\frac{\frac{d}{dh}(\ln(1+h))}{\frac{d}{dh}(h)}\\
 & = & \lim_{h\to0}\frac{1}{1+h}\\
 & = & 1.
\end{eqnarray*}
Dengan demikian, $L=1$, dan karena fungsi eksponensial $e^{x}$ kontinu,
\begin{eqnarray*}
e=e^{L}=e^{\lim_{h\to0}\frac{\ln(1+h)}{h}}=\lim_{h\to0}e^{\frac{\ln(1+h)}{h}} & = & \lim_{h\to0}e^{\ln\left[(1+h)^{1/h}\right]}\\
 & = & \lim_{h\to0}(1+h)^{1/h}.
\end{eqnarray*}
\item [{$\tilde{e}_{4}(h)$:}] \label{eTilde4}Kita menggunakan rumus \ref{eq:richardsonBasis}
dengan $\alpha=\frac{1}{2}$, $m=4$, dan $n=5$ untuk memperoleh
\begin{eqnarray*}
\tilde{e}_{4}(h) & = & \frac{16\tilde{e}_{3}\left(\frac{h}{2}\right)-\tilde{e}_{3}(h)}{15}\\
 & = & e+O(h^{5}).
\end{eqnarray*}
Dengan menerapkan rumus ini pada $\tilde{e}_{3}(0.01)$ dan $\tilde{e}_{3}(0.005)$,
kita memperoleh
\begin{eqnarray*}
\tilde{e}_{4}(0.01) & = & \frac{16(2.718281828448482)-2.718281828281785}{15}\\
 & = & 2.718281828459595,
\end{eqnarray*}
nilai yang akurat hingga 13 angka signifikan!
\item [{rasio~galat:}] \label{errorRatio}Kita mengharapkan $\frac{\left|\int_{0}^{1}\sin x\,dx-T_{10}\right|}{\left|\int_{0}^{1}\sin x\,dx-T_{15}\right|}$
bernilai kira-kira $1.5^{2}=2.25$ karena 15 (jumlah subselang yang
digunakan dalam hampiran penyebut) adalah $1.5$ kali 10 (jumlah subselang
yang digunakan dalam hampiran pembilang).
\item [{$T_{5,3}$~dan~galat~mutlaknya:}] \label{t35}$\frac{\left|\int_{0}^{1}\sin x\,dx-T_{5,2}\right|}{\left|\int_{0}^{1}\sin x\,dx-T_{10,2}\right|}\approx\frac{1.522(10)^{-11}}{2.374(10)^{-13}}\approx64$
sehingga 
\begin{eqnarray*}
T_{5,3} & = & \frac{64T_{10,2}-T_{5,2}}{63}\approx.4596976941318606\\
\left|\int_{0}^{1}\sin x\,dx-T_{5,3}\right| & \approx & 4(10)^{-16}
\end{eqnarray*}
\end{description}

